% Part: normal-modal-logic
% Chapter: axioms-systems
% Section: proofs-in-K
%READERNOTE{SOL6-A156: दहा ओळींच्या K-सिद्धतेत ओळ 6 मधील implication चे बाह्य कोष्ठक उघडले आणि ओळी 8 व 9 मधील प्रत्येकी एक जादा बंद कोष्ठक काढले. ओळी 3 व 5 चे K-दृष्टांत आणि सर्व inference references आधीच योग्य होते; ते जपले.}

\documentclass[../../../include/open-logic-section]{subfiles}

\begin{document}

\olfileid{nml}{prf}{prk}

\olsection{\Log{K} मधील सिद्धता}

लघुतम मोडल प्रणालीतील सिद्धतेचा सराव म्हणून
\olref[syn][sch]{tab:valid-invalidSchemas} च्या
डाव्या बाजूची वैध !!{formula}s सर्व
\Log{K} मध्ये सिद्ध करता येतात, हे दाखवू.

\begin{prop}
  $\Log{K} \Proves \Box!A \lif \Box (!B\lif !A)$
\end{prop}

\begin{proof}
  \begin{derivation}
    1. & $!A \lif (!B \lif !A)$ & \Taut \\
    2. & $\Box(!A \lif (!B \lif !A))$ & \Nec, 1 \\
    3. & $\Box(!A \lif (!B \lif !A)) \lif
    (\Box!A \lif \Box(!B \lif !A))$ & \Ax{K}\\
    4. & $\Box!A \lif \Box (!B\lif !A)$ & \MP, 2, 3
  \end{derivation}
\end{proof}

\begin{prop}
  $\Log{K} \Proves \Box(!A \land !B) \lif (\Box !A \land \Box!B)$
\end{prop}

\begin{proof}
\begin{derivation}
  1. & $(!A \land !B) \lif !A$ & \Taut \\
  2. & $\Box((!A \land !B) \lif !A)$ & \Nec \\
  3. & $\Box((!A \land !B) \lif !A) \lif (\Box(!A \land !B) \lif \Box!A)$
  & \Ax{K} \\
  4. & $\Box(!A \land !B) \lif \Box!A$ & \MP, 2, 3 \\
  5. & $(!A \land !B) \lif !B$ & \Taut \\
  6. & $\Box((!A \land !B) \lif !B)$ & \Nec \\
  7. & $\Box((!A \land !B) \lif !B) \lif (\Box(!A \land !B) \lif \Box!B)$
  & \Ax{K} \\
  8. & $\Box(!A \land !B) \lif \Box!B$ & \MP, 6, 7 \\
  9. & $(\Box(!A \land !B) \lif \Box!A) \lif{}$ \\
  & \qquad $((\Box(!A \land !B) \lif \Box!B) \lif{}$ \\
  & \qquad $(\Box(!A \land !B) \lif (\Box !A \land \Box!B)))$ & \Taut\\
  10. &  $(\Box(!A \land !B) \lif \Box!B) \lif{}$ \\
  & \qquad $(\Box(!A \land !B) \lif (\Box !A \land \Box!B))$ & \MP, 4, 9\\
  11. & $\Box(!A \land !B) \lif (\Box !A \land \Box!B)$ & \MP, 8, 10.
\end{derivation}
ओळ~$9$ वरील !!{formula} पुढील सर्वतः-सत्य सूत्राचा दृष्टांत आहे:
\[
(p \lif q) \lif ((p \lif r) \lif (p \lif (q \land r))).
\]
\end{proof}

\begin{prop}
  $\Log{K}\Proves (\Box!A \land \Box!B) \lif \Box (!A \land !B)$
\end{prop}

\begin{proof}
  \begin{derivation}
    1. & $!A \lif (!B \lif (!A \land !B))$ & \Taut \\
    2. & $\Box(!A \lif (!B \lif (!A \land !B)))$ & \Nec, 1 \\
    3. & $\Box(!A \lif (!B \lif (!A \land !B))) \lif (\Box!A \lif \Box(!B \lif (!A \land !B)))$ & \Ax{K}\\
    4. & $\Box!A \lif \Box(!B \lif (!A \land !B))$ & \MP, 2, 3\\
    5. & $\Box(!B \lif (!A \land !B)) \lif (\Box!B \lif \Box(!A \land !B))$ & \Ax{K} \\
    6. & $(\Box!A \lif \Box(!B \lif (!A \land !B))) \lif {}$ \\
    & \qquad $((\Box(!B \lif (!A \land !B)) \lif (\Box!B \lif \Box(!A \land !B))) \lif {}$\\
    & \qquad $(\Box!A \lif (\Box !B \lif \Box(!A \land !B))))$ & \Taut\\
    7. & $(\Box(!B \lif (!A \land !B)) \lif (\Box!B \lif \Box(!A \land !B))) \lif {}$\\
    & \qquad $(\Box!A \lif (\Box !B \lif \Box(!A \land !B)))$ & \MP, 4, 6\\
    8. & $\Box!A \lif (\Box !B \lif \Box(!A \land !B))$ & \MP, 5, 7\\
    9. & $(\Box!A \lif (\Box !B \lif \Box(!A \land !B))) \lif {}$\\
    & \qquad $((\Box!A \land \Box!B) \lif \Box(!A \land !B))$ & \Taut\\
    10. & $(\Box!A \land \Box!B) \lif \Box(!A \land !B)$ & \MP, 8, 9
  \end{derivation}
  ओळी $6$ आणि $9$ वरील !!{formula}s अनुक्रमे पुढील
  सर्वतः-सत्य सूत्रांचे दृष्टांत आहेत:
  \begin{align*}
    (p \lif q) & \lif ((q \lif r) \lif (p \lif r)) \\
    (p \lif (q \lif r)) & \lif ((p \land q) \lif r)
  \end{align*}
\end{proof}

\begin{prop}
  \iftag{prvBox}{$\Log{K} \Proves \lnot\Box p \lif \Diamond \lnot
    p$}{$\Log{K} \Proves \Box\lnot p \lif \lnot\Diamond p$}
\end{prop}

\begin{proof}
  \iftag{prvBox}{\iftag{prvDiamond}{% Box and Diamond are primitive
  \begin{derivation}
    1. & $\Diamond \lnot p \liff \lnot\Box\lnot\lnot p$ & \Dual\\
    2. & $(\Diamond \lnot p \liff \lnot\Box\lnot\lnot p) \lif {}$ \\
    & \qquad $(\lnot \Box \lnot\lnot p \lif \Diamond \lnot p)$& \Taut\\
    3. & $\lnot \Box \lnot\lnot p \lif \Diamond \lnot p$ & \MP, 1, 2\\
    4. & $\lnot\lnot p \lif p$ & \Taut\\
    5. & $\Box(\lnot\lnot p \lif p)$ & \Nec, 4\\
    6. & $\Box(\lnot\lnot p \lif p) \lif (\Box \lnot\lnot p \lif \Box p)$ & \Ax{K}\\
    7. & $(\Box \lnot\lnot p \lif \Box p)$ & \MP, 5, 6\\
    8. & $(\Box \lnot\lnot p \lif \Box p) \lif (\lnot \Box p \lif \lnot\Box\lnot\lnot p)$ & \Taut\\
    9. & $\lnot \Box p \lif \lnot\Box\lnot\lnot p$ & \MP, 7, 8\\
    10. & $(\lnot \Box p \lif \lnot\Box\lnot\lnot p) \lif {}$ \\
    & \qquad $((\lnot \Box \lnot\lnot p \lif \Diamond \lnot p) \lif (\lnot \Box p \lif \Diamond\lnot p))$ & \Taut\\
    11. & $(\lnot \Box \lnot\lnot p \lif \Diamond \lnot p) \lif (\lnot \Box p \lif \Diamond\lnot p)$ & \MP, 9, 10\\
    12. & $\lnot\Box p \lif \Diamond \lnot p$ & \MP, 3, 11\\
  \end{derivation}
  ओळी $8$ आणि $10$ वरील !!{formula}s अनुक्रमे पुढील
  सर्वतः-सत्य सूत्रांचे दृष्टांत आहेत:
  \begin{align*}
    & (p \lif q) \lif (\lnot q \lif \lnot p) \\
    & (p \lif q) \lif ((q \lif r) \lif (p \lif r)).
  \end{align*}
  }{% Diamond is defined
  \begin{derivation}
    1. & $\lnot\lnot p \lif p$ & \Taut\\
    2. & $\Box(\lnot\lnot p \lif p)$ & \Nec, 1\\
    3. & $\Box(\lnot\lnot p \lif p) \lif (\Box \lnot\lnot p \lif \Box p)$ & \Ax{K}\\
    4. & $(\Box \lnot\lnot p \lif \Box p)$ & \MP, 2, 3\\
    5. & $(\Box \lnot\lnot p \lif \Box p) \lif (\lnot \Box p \lif \lnot\Box\lnot\lnot p)$ & \Taut\\
    6. & $\lnot \Box p \lif \lnot\Box\lnot\lnot p$ & \MP, 4, 5\\
    & $\lnot\Box p \lif \Diamond \lnot p$
  \end{derivation}
  अंतिम ओळ $\lnot \Box p \lif \Diamond\lnot p$ अशी आहे, कारण
  $\Diamond$ ची व्याख्या $\lnot\Box\lnot$ अशी आहे.
  ओळ~$5$ वरील !!{formula} पुढील सर्वतः-सत्य सूत्राचा दृष्टांत आहे:
  \[
    (p \lif q) \lif (\lnot q \lif \lnot p).
  \]
  }}{% Box is defined
  \begin{derivation}
    1. & $\Diamond  p \liff \lnot\Box\lnot p$ & \Dual\\
    2. & $(\Diamond p \liff \lnot\Box\lnot p) \lif
    (\Diamond p \lif \lnot \Box \lnot p)$& \Taut\\
    3. & $\Diamond p \lif \lnot \Box \lnot p$ & \MP, 1, 2\\
    4. & $(\Diamond p \lif \lnot \Box \lnot p) \lif (\Box\lnot p \lif \lnot\Diamond p)$ & \Taut\\
    5. & $\Box\lnot p \lif \lnot\Diamond p$ & \MP, 3, 4
  \end{derivation}
  ओळ~$4$ वरील !!{formula} पुढील सर्वतः-सत्य सूत्राचा
  दृष्टांत आहे:
  \begin{align*}
    & (p \lif \lnot q) \lif (q \lif \lnot p).
  \end{align*}
}
\end{proof}

\begin{prob}
  पुढील !!{formula}s साठी $\Log{K}$ मध्ये
  !!{derivation}s शोधा:
  \begin{enumerate}
    \item $\Box \lnot p \lif \Box(p \lif q)$
    \item $(\Box p \lor \Box q) \lif \Box(p \lor q)$
    \item $\Diamond p \lif \Diamond(p \lor q)$
  \end{enumerate}
\end{prob}

\end{document}
